\section{Discussion}\label{sec:discussion}
The first question concerns the evolution of the filtered genus. The measured curves change modestly across the four snapshots, while foreground removal changes their amplitudes and shapes substantially. Volume-fraction thresholds remove sensitivity to a uniform rescaling of density contrast, but do not remove evolution from changing morphology or tracer selection. The amplitude changes reported here are descriptive; their cross-redshift significance has not been calibrated.

The second question concerns the response to geometry. The tested incorrect mappings leave residuals after subtraction of the simulation-cosmology template, including in foreground-cut fields. This supports further investigation of redshift-dependent topology, but the numerical tests show that a response of the pipeline is not automatically a stable physical constraint. The simulation-cosmology zero is imposed by construction, and neither the template nor the regularized covariance has been validated with independent realizations. Internal resampling also omits modes larger than the available volume \citep{LacasaKunz2017InternalSSC}.

The third question concerns the meaning of the parameter areas. Mesh and mask discontinuities, empirical smoothing and finite-domain normalization materially affect them. Spatial convergence and independent statistical calibration are needed before interpreting them as parameter precision \citep{Kim2014GenusSystematics,Appleby2017,ApplebyBOSS2022}. The idealized cuts probe mode-loss geometries rather than an instrument-specific foreground model; realistic observations additionally require signal-loss calibration and survey selection modelling \citep{Cunnington2020ForegroundPNG,Cunnington2023ForegroundTransfer}. The present study identifies both a measurable geometric response and the numerical tests needed to develop it into a quantitative probe.
